\documentclass{article}
\usepackage{amsmath, amssymb, amsthm}
\usepackage{hyperref}
\usepackage{geometry}
\geometry{margin=1in}

\title{Erd\H{o}s Problem \#346}
\date{Last updated: 2025-08-31}
\author{Source: erdosproblems.com}

\begin{document}
\maketitle

\noindent\textbf{Status:} OPEN \quad \textbf{No prize}

\noindent\textbf{Tags:} number theory, complete sequences

\medskip
\noindent\textbf{Problem Statement:}

Let $A=\{1\leq a_1< a_2<\cdots\}$ be a set of integers such that
{UL}
{LI} $A\backslash B$ is complete for any finite subset $B$ and {/LI}
{LI} $A\backslash B$ is not complete for any infinite subset $B$.{/LI}
{/UL}
(Here 'complete' means all sufficiently large integers can be written as a sum of distinct members of the sequence.)

Is it true that if $a_{n+1}/a_n \geq 1+\epsilon$ for some $\epsilon>0$ and all $n$ then\[\lim_n \frac{a_{n+1}}{a_n}=\frac{1+\sqrt{5}}{2}?\]

\bigskip
\noindent\textbf{Remarks:}

Graham \cite{Gr64d} has shown that the sequence $a_n=F_n-(-1)^{n}$, where $F_n$ is the $n$th Fibonacci number, has these properties. Erd\H{o}s and Graham \cite{ErGr80} remark that it is easy to see that if $a_{n+1}/a_n>\frac{1+\sqrt{5}}{2}$ then the second property is automatically satisfied, and that it is not hard to construct very irregular sequences satisfying both properties.

\bigskip
\noindent\textbf{References:}

[ErGr80] Erd\H{o}s, P. and Graham, R., Old and new problems and results in combinatorial number theory. Monographies de L'Enseignement Mathematique (1980).
 
 [Gr64d] Graham, R. L., A property of Fibonacci numbers. Fibonacci Quart. (1964), 1-10.

\bigskip
\noindent\small{Source: \url{https://www.erdosproblems.com/346}}
\end{document}
